\chapter{Inception and the Executable Pipeline (Python 1.0--1.6)}

Python's surface is famously simple; its execution model is not. To master the language
you must first hold one mental picture firmly: \textbf{source text is not run---it is
compiled, ahead of time, into an immutable \texttt{PyCodeObject}, which a stack-based
virtual machine then interprets.} This chapter follows a single line of source from the
tokenizer to the evaluation loop, establishing the vocabulary---names, bindings, scopes,
cells, code objects, frames, opcodes---that every later chapter assumes. We anchor in
Python's 1.x origins because the design decisions made then (names as bindings,
everything-is-an-object, a static scope-analysis pass before execution) still govern how
CPython 3.13 behaves today.

\section*{Section Index}
\begin{itemize}
  \item \textbf{1.1} Inception and the core design philosophy
  \item \textbf{1.2} Names are bindings, not storage: the object model and the senior-engineer contrast
  \item \textbf{1.3} The execution pipeline, end to end
  \item \textbf{1.4} Parser evolution: LL(1) to PEG (PEP 617)
  \item \textbf{1.5} The symbol table and static scope classification
  \item \textbf{1.6} Closures and cell objects
  \item \textbf{1.7} \texttt{PyCodeObject}: the compiled blueprint
  \item \textbf{1.8} The virtual machine: frames and the evaluation loop
  \item \textbf{1.9} Performance model and anti-patterns
  \item \textbf{1.10} Summary and cross-references
\end{itemize}

\section{Inception and the core design philosophy}

\textbf{Why this exists.} Python was conceived in December 1989 by Guido van Rossum at CWI
(Centrum Wiskunde \& Informatica) in the Netherlands, first released publicly as version
0.9.0 in February 1991, and reached 1.0 in January 1994. It was a deliberate successor to
the \textbf{ABC language}---elegant for teaching, but crippled in practice by three
architectural decisions Python set out to reverse:

\begin{enumerate}
  \item \textbf{A monolithic, unextensible runtime.} ABC could not load native libraries or
    be driven from C. Python was built around a C-extension boundary: the interpreter
    \emph{is} a C library, and any data structure or algorithm can be dropped to C without
    changing the Python-level interface. This is why NumPy, the \texttt{re} engine, and
    \texttt{asyncio}'s selectors can exist at all.
  \item \textbf{A global, filesystem-backed state model.} ABC persisted variables in a
    global database, making it slow and impossible to compose. Python uses ordinary
    in-memory namespaces (dictionaries) with lexical scoping.
  \item \textbf{A closed ecosystem.} ABC could not reach the operating system. Python
    exposed system calls, signals, and process control as first-class library surface from
    its earliest versions.
\end{enumerate}

Three design commitments fell out of these reversals and never changed:

\begin{itemize}
  \item \textbf{Significant whitespace.} CPython's tokenizer emits explicit \texttt{INDENT}
    and \texttt{DEDENT} tokens; block structure is part of the grammar, not a convention.
    The \emph{visual} structure a reader sees is, by construction, the \emph{execution}
    structure the compiler sees---the two cannot drift apart, as they can in brace
    languages.
  \item \textbf{A unified object model.} Every value---an integer, a function, a class, a
    module, a stack frame---is a heap-allocated object reachable through a
    \texttt{PyObject *}. There is no separate world of ``primitives'' as in Java or C\#.
    This uniformity is what makes introspection total rather than partial.
  \item \textbf{Names as bindings.} \texttt{x = 5} does not create a typed memory slot named
    \texttt{x}; it binds the \emph{name} \texttt{x}, in some namespace, to the object
    \texttt{5}. This is the most consequential decision in the language, and §1.2 unpacks it.
\end{itemize}

\textbf{What the interpreter actually does at startup.} Running \texttt{python script.py}
initializes a process-wide runtime (\texttt{PyInterpreterState}) and at least one thread
state (\texttt{PyThreadState}), imports the \texttt{builtins}, \texttt{sys}, and
\texttt{\_\_main\_\_} modules, then compiles and executes \texttt{script.py} in the
\texttt{\_\_main\_\_} namespace. ``Running a script'' is therefore ``compile a module body to
a code object, then evaluate it in a fresh frame.''

\section{Names are bindings, not storage: the object model and the senior-engineer contrast}

\textbf{The model.} A Python name lives in a namespace---a mapping from \texttt{str} to
\texttt{PyObject *}. Assignment rebinds the name; it never copies or mutates the object the
name previously referred to. Every object carries, at minimum, a \textbf{reference count}
and a \textbf{type pointer} (Chapter 2 dissects the \texttt{PyObject} header). Assignment is
a constant two-step: bind the name, and increment the target object's refcount.

\textbf{The senior-engineer contrast.} This is where intuition from systems languages
misleads:

\begin{itemize}
  \item \textbf{What a variable is.} C/C++: typed storage at an address. Java: a typed slot
    holding a primitive or a reference. Python: a name bound to an object reference; the
    name has no type.
  \item \textbf{\texttt{a = b}.} C/C++: copies bytes per the declared type. Java: copies a
    primitive or a reference. Python: binds \texttt{a} to the same object \texttt{b} refers
    to; no copy.
  \item \textbf{Type.} C/C++ and Java: belongs to the variable. Python: belongs to the
    \emph{object}, not the name.
  \item \textbf{Lifetime.} C/C++: scope / \texttt{new}--\texttt{delete} / RAII. Java:
    tracing GC. Python: refcount $\rightarrow$ immediate free, plus a cyclic collector
    (Chapter 3).
\end{itemize}

\begin{lstlisting}[language=Python, caption={Assignment binds; it does not copy. Mutating through one name is visible through every name bound to the same object.}]
a = [1, 2, 3]
b = a            # b is NOT a copy; a and b name the same list object
b.append(4)
print(a)         # [1, 2, 3, 4]
print(a is b)    # True  -> identity, not equality
\end{lstlisting}

Verified output (CPython 3.13.5):

\begin{lstlisting}[caption={Output}]
[1, 2, 3, 4]
True
\end{lstlisting}

\texttt{is} compares \textbf{identity} (are these the same object?), \texttt{==} compares
\textbf{value}. The distinction is meaningless in C, where \texttt{a == b} on pointers
\emph{is} identity; in Python the two are separate operators because a name and an object are
separate things. Fix the model: \textbf{a Python program is a graph of objects, and names
are merely labelled edges into it.}

\section{The execution pipeline, end to end}

\textbf{Why this exists.} Python feels interpreted, but nothing executes until it has been
compiled to bytecode. The pipeline tells you \emph{where} each kind of error is produced (a
\texttt{SyntaxError} cannot occur at runtime; a \texttt{NameError} cannot occur at compile
time) and \emph{why} constructs are fast or slow.

\begin{lstlisting}[caption={From text to result}]
[ source text ]
   | tokenizer (Parser/tokenizer.c): token stream, incl. INDENT/DEDENT
   v
[ token stream ]
   | parser (PEG since 3.9): builds the Abstract Syntax Tree (AST)
   v
[   AST   ]   ( the ast module exposes this tree )
   | symbol-table pass (Python/symtable.c): classifies every name into a scope
   v
[ symbol table ]
   | compiler (Python/compile.c): emits bytecode + metadata
   v
[ PyCodeObject ]   ( immutable: bytecode, constants, names, flags )
   | evaluation loop (Python/ceval.c, _PyEval_EvalFrameDefault): runs per frame
   v
[  result  ]
\end{lstlisting}

The first four stages are \textbf{ahead-of-time compilation}; only the last is
``interpretation.'' Each boundary is observable from the standard library: \texttt{tokenize}
for stage 1, \texttt{ast} for stages 2--3, \texttt{symtable} for stage 4, \texttt{compile()}
and \texttt{dis} for stage 5's input.

\section{Parser evolution: LL(1) to PEG (PEP 617)}

\textbf{Why this exists.} For most of its life CPython used a hand-tunable \textbf{LL(1)}
parser---\textbf{L}eft-to-right scan, \textbf{L}eftmost derivation, \textbf{1} token of
lookahead. LL(1) is fast and simple but structurally limited. \textbf{PEP 617} replaced it
with a \textbf{PEG} parser in Python 3.9.

\textbf{The LL(1) limitation: left recursion.} A top-down LL(1) parser cannot directly
handle a \textbf{left-recursive} rule---one whose right-hand side begins with the
non-terminal itself:
\[ A \rightarrow A\,\alpha \;\mid\; \beta \]
Asked to expand $A$, the parser would expand $A$ again before consuming any input,
recursing forever:
\[ A \rightarrow A\alpha \rightarrow A\alpha\alpha \rightarrow A\alpha\alpha\alpha \rightarrow \cdots \]
The classic workaround mechanically rewrites the rule into an equivalent right-recursive
pair:
\[ A \rightarrow \beta A' \qquad A' \rightarrow \alpha A' \;\mid\; \varepsilon \]
This is equivalent as a \emph{language} but distorts the grammar: the natural
left-associative shape of \texttt{a - b - c} must be reconstructed afterward.

\textbf{The PEG parser} is also top-down, but adds \textbf{ordered choice} (first match
wins), \textbf{unlimited lookahead} via backtracking, \textbf{packrat memoization} (caching
each rule/position result, which both bounds backtracking to linear time and supports left
recursion directly), and \textbf{direct AST construction}, skipping the intermediate Concrete
Syntax Tree.

\begin{lstlisting}[caption={Parser pipelines, old and new}]
LL(1) pipeline (pre-3.9):
[source] -> [tokenizer] -> [Concrete Syntax Tree] -> [AST builder] -> [AST]

PEG pipeline (3.9+, PEP 617):
[source] -> [tokenizer] -> [PEG parser (memoized)] -> [AST]
\end{lstlisting}

\textbf{Old-vs-new, why it mattered.} The \textbf{structural pattern matching} of
\texttt{match}/\texttt{case} (PEP 634, Chapter 15) is expressible largely \emph{because} the
grammar is no longer bound by LL(1). The parser is an implementation detail---the
\emph{language} is defined by the grammar in the reference (covered in Vol X)---but the
switch enabled strictly more expressive syntax and better diagnostics.

\section{The symbol table and static scope classification}

\textbf{Why this exists.} Before emitting a single opcode, the compiler walks the AST to
decide, \textbf{statically}, what every name \emph{means}: local, parameter, global,
builtin, or free variable. This classification lets CPython use array-indexed local access
(\texttt{LOAD\_FAST}) instead of a dictionary lookup per variable---the single largest reason
Python functions are not even slower than they are.

CPython classifies each binding (\texttt{Python/symtable.c}) as: \textbf{local} (bound in the
current function); \textbf{global, explicit} (\texttt{global x}); \textbf{global, implicit}
(referenced but never bound; resolved at runtime against module globals, then builtins); or
\textbf{enclosing/free} (bound in an outer function and read, or with \texttt{nonlocal}
rebound, by a nested function).

\begin{lstlisting}[language=Python, caption={Inspect CPython's static scope analysis without running the code.}]
import symtable

src = """
def outer_func(x):
    z = 10
    def inner_func():
        nonlocal z
        return x + z
    return inner_func
"""

table = symtable.symtable(src, filename="<string>", compile_type="exec")
outer = table.lookup("outer_func").get_namespace()
print("outer identifiers:", sorted(outer.get_identifiers()))
print("'z' is_local in outer:", outer.lookup("z").is_local())
print("'x' is_parameter in outer:", outer.lookup("x").is_parameter())

inner = outer.lookup("inner_func").get_namespace()
print("inner free vars (closure):", inner.get_frees())
print("'z' is_free in inner:", inner.lookup("z").is_free())
\end{lstlisting}

Verified output (CPython 3.13.5):

\begin{lstlisting}[caption={Output}]
outer identifiers: ['inner_func', 'x', 'z']
'z' is_local in outer: True
'x' is_parameter in outer: True
inner free vars (closure): ('z', 'x')
'z' is_free in inner: True
\end{lstlisting}

\textbf{Correctness note.} Earlier drafts called \texttt{Symbol.is\_cell()}. That method does
not exist in the \texttt{symtable} API (verified on 3.13.5; it raises
\texttt{AttributeError}). The supported predicates are \texttt{is\_local}, \texttt{is\_global},
\texttt{is\_free}, \texttt{is\_parameter}, \texttt{is\_namespace}, \texttt{is\_nonlocal},
\texttt{is\_declared\_global}, \texttt{is\_assigned}, \texttt{is\_referenced},
\texttt{is\_imported}, \texttt{is\_annotated}. To see which locals became \textbf{cells}, read
the compiled code object's \texttt{co\_cellvars} (§1.6).

\section{Closures and cell objects}

\textbf{Why this exists.} A nested function may outlive the call that created it, yet still
read the enclosing function's locals. Those locals cannot live in the outer frame's
fast-locals array, because that array dies when the frame is popped. CPython promotes a
captured local to a \textbf{cell}---a tiny heap object shared by reference between the
enclosing and nested functions.

\begin{lstlisting}[language=C, caption={Illustrative; Objects/cellobject.c. A cell is a one-slot indirection on the heap.}]
typedef struct {
    PyObject_HEAD
    PyObject *ob_ref;   /* the captured object, shared by closure and enclosing scope */
} PyCellObject;
\end{lstlisting}

The compiler records, on each code object, \texttt{co\_cellvars} (locals of \emph{this}
function captured by a nested function) and \texttt{co\_freevars} (names \emph{this} function
reads from an enclosing scope).

\begin{lstlisting}[language=Python, caption={The code-object view of cells and free variables (authoritative).}]
src = """
def outer_func(x):
    z = 10
    def inner_func():
        nonlocal z
        return x + z
    return inner_func
"""
top = compile(src, "<string>", "exec")
(outer_code,) = [c for c in top.co_consts if getattr(c, "co_name", None) == "outer_func"]
(inner_code,) = [c for c in outer_code.co_consts if getattr(c, "co_name", None) == "inner_func"]
print("outer_func co_cellvars:", outer_code.co_cellvars)
print("inner_func co_freevars:", inner_code.co_freevars)
\end{lstlisting}

Verified output (CPython 3.13.5):

\begin{lstlisting}[caption={Output}]
outer_func co_cellvars: ('x', 'z')
inner_func co_freevars: ('x', 'z')
\end{lstlisting}

Both \texttt{x} and \texttt{z} are cells in \texttt{outer\_func} and free variables of
\texttt{inner\_func}. The shared cell is why a closure sees \emph{live} updates to a captured
variable, not a snapshot---the source of the ``late-binding closure in a loop'' bug dissected
in Chapter 3.

\section{\texttt{PyCodeObject}: the compiled blueprint}

\textbf{Why this exists.} Compilation produces a \texttt{PyCodeObject}: an \textbf{immutable},
shareable description of an executable unit (module body, function, \texttt{lambda},
comprehension, class body). Immutability lets it be cached, shared across calls and threads,
and reasoned about statically. A \emph{function} is a thin, mutable wrapper
(\texttt{\_\_defaults\_\_}, \texttt{\_\_closure\_\_}, \texttt{\_\_globals\_\_}) around an
immutable code object.

\begin{lstlisting}[language=C, caption={Illustrative; fields from Include/cpython/code.h. The struct changed across versions; these co\_* fields are stable and observable from Python.}]
struct PyCodeObject {
    PyObject_HEAD
    int co_argcount;        /* positional parameters */
    int co_posonlyargcount; /* positional-only parameters (PEP 570, 3.8+) */
    int co_kwonlyargcount;  /* keyword-only parameters */
    int co_nlocals;         /* number of local variables */
    int co_stacksize;       /* max evaluation-stack depth the VM must reserve */
    int co_flags;           /* compiler flags (CO_OPTIMIZED, CO_NEWLOCALS, CO_GENERATOR...) */
    PyObject *co_consts;    /* tuple of literal constants */
    PyObject *co_names;     /* tuple of global/attribute names */
    PyObject *co_varnames;  /* tuple of local variable names */
    PyObject *co_freevars;  /* tuple of free (captured) names */
    PyObject *co_cellvars;  /* tuple of local names boxed into cells */
};
\end{lstlisting}

\begin{lstlisting}[language=Python, caption={Read the compiled blueprint of a live function.}]
def example_fn(a, b):
    local_val = 42
    return a + b + local_val

co = example_fn.__code__
print("co_varnames:", co.co_varnames)
print("co_consts:", co.co_consts)
print("co_argcount:", co.co_argcount, "co_nlocals:", co.co_nlocals,
      "co_stacksize:", co.co_stacksize)
print("co_flags:", bin(co.co_flags))
\end{lstlisting}

Verified output (CPython 3.13.5):

\begin{lstlisting}[caption={Output}]
co_varnames: ('a', 'b', 'local_val')
co_consts: (None, 42)
co_argcount: 2 co_nlocals: 3 co_stacksize: 2
co_flags: 0b11
\end{lstlisting}

\begin{itemize}
  \item \texttt{co\_consts} contains \texttt{None} even though the source never mentions
    it---the implicit return value's constant. The literal \texttt{42} is interned into the
    pool; bytecode refers to it by index.
  \item \texttt{co\_flags == 0b11} sets \texttt{CO\_OPTIMIZED | CO\_NEWLOCALS}: ``this is a
    real function---use fast array-indexed locals and a fresh locals namespace per call.'' A
    module or class body does \emph{not} set these, which is why module-level code uses
    dictionary-based \texttt{LOAD\_NAME} and is slower (§1.9).
  \item \textbf{\texttt{EXTENDED\_ARG}.} An opcode's inline argument is one byte. To reference
    the 300th constant, the compiler prefixes \texttt{EXTENDED\_ARG} to supply the high-order
    bits; the VM shifts and combines them, allowing arguments up to 32 bits wide.
\end{itemize}

\section{The virtual machine: frames and the evaluation loop}

\textbf{Why this exists.} CPython is a \textbf{stack-based virtual machine}. Each call
creates a \textbf{frame} holding the code object, the locals, and an \textbf{evaluation
stack}. A single monolithic C function, \texttt{\_PyEval\_EvalFrameDefault} in
\texttt{Python/ceval.c}, dispatches opcodes in a loop until the frame returns.

\begin{lstlisting}[language=C, caption={Illustrative dispatch loop, Python/ceval.c. The modern loop uses computed-goto threaded dispatch and inline caches, but the skeleton is this.}]
PyObject *_PyEval_EvalFrameDefault(PyThreadState *tstate,
                                   _PyInterpreterFrame *frame, int throwflag) {
    for (;;) {
        _Py_CODEUNIT word = *next_instr++;
        switch (word.op.code) {
            case LOAD_FAST: {                 /* push a local by array index  */
                PyObject *v = GETLOCAL(word.op.arg);
                Py_INCREF(v); PUSH(v); DISPATCH();
            }
            case BINARY_OP: {                 /* pop two, push their combination */
                PyObject *r = POP(), *l = POP();
                PyObject *res = PyNumber_BinaryOp(l, r, word.op.arg);
                Py_DECREF(l); Py_DECREF(r); PUSH(res); DISPATCH();
            }
            case RETURN_VALUE:
                return POP();
        }
    }
}
\end{lstlisting}

Do not trust the textbook picture---read what \emph{this} interpreter emits. The \texttt{dis}
module disassembles a code object into the real instruction stream:

\begin{lstlisting}[language=Python, caption={The actual bytecode CPython 3.13 runs for example\_fn.}]
def example_fn(a, b):
    local_val = 42
    return a + b + local_val

import dis
dis.dis(example_fn)
\end{lstlisting}

Verified output (CPython 3.13.5):

\begin{lstlisting}[caption={Output}]
  1           RESUME                   0

  2           LOAD_CONST               1 (42)
              STORE_FAST               2 (local_val)

  3           LOAD_FAST_LOAD_FAST      1 (a, b)
              BINARY_OP                0 (+)
              LOAD_FAST                2 (local_val)
              BINARY_OP                0 (+)
              RETURN_VALUE
\end{lstlisting}

This listing demolishes several outdated mental models:

\begin{itemize}
  \item \textbf{\texttt{RESUME 0}} opens every code object since 3.11: a no-op hook used by
    the specializing adaptive interpreter and by debuggers/profilers. Its presence marks
    modern bytecode.
  \item \textbf{\texttt{BINARY\_OP 0 (+)}}, not \texttt{BINARY\_ADD}. Python 3.11 collapsed
    the per-operator binary opcodes into a single \texttt{BINARY\_OP} whose argument selects
    the operation. Material showing \texttt{BINARY\_ADD} predates 3.11.
  \item \textbf{\texttt{LOAD\_FAST\_LOAD\_FAST 1 (a, b)}} is a \textbf{superinstruction}: the
    compiler fused two adjacent \texttt{LOAD\_FAST}s to cut dispatch overhead. The VM is no
    longer one opcode per source operation.
  \item There is \textbf{no trailing \texttt{LOAD\_CONST None} / \texttt{RETURN\_VALUE}}
    pair: the explicit \texttt{return} means the compiler omits the implicit
    \texttt{return None}.
\end{itemize}

\begin{lstlisting}[caption={Evaluation-stack trace for a + b}]
  LOAD_FAST_LOAD_FAST(a,b)   BINARY_OP(+)        ...        RETURN_VALUE
  +---------+                +---------+                    +---------+
  |    b    |                |  a + b  |                    | (empty) |
  +---------+                +---------+        -->         +---------+
  |    a    |       -->      | (empty) |                    | (empty) |
  +---------+                +---------+                    +---------+
\end{lstlisting}

\textbf{3.11+ reality: the specializing adaptive interpreter.} Modern CPython does not
execute a fixed opcode stream forever. As code runs hot, generic opcodes are
\textbf{specialized} in place into faster variants, with \textbf{inline caches} stored in the
code units themselves. This is \emph{not} the JIT (PEP 744, Vol VII)---it is interpretation
that rewrites itself. Covered in depth in Vol VI, Ch 17.

\section{Performance model and anti-patterns}

\textbf{Local access is array indexing; global/builtin access is dictionary lookup.} Inside
an optimized function, a local read is \texttt{LOAD\_FAST} (an index into a C array). A global
or builtin read is \texttt{LOAD\_GLOBAL}, which probes the module \texttt{\_\_dict\_\_} and
then \texttt{builtins}. The classic hot-loop optimization---binding a global or method to a
local before the loop---exploits exactly this difference:

\begin{lstlisting}[language=Python, caption={Hoisting a global lookup to a local turns N dict probes into N array reads.}]
import math

def slow(xs):
    return [math.sqrt(x) for x in xs]      # LOAD_GLOBAL math; LOAD_ATTR sqrt -- each item

def fast(xs):
    sqrt = math.sqrt                       # one lookup, hoisted
    return [sqrt(x) for x in xs]           # LOAD_FAST sqrt -- each item
\end{lstlisting}

Both are correct; \texttt{fast} pays the lookup once. This is a constant-factor win, not an
algorithmic one---reach for it only in genuine hot loops (see the profiling discipline in
Vol IX).

\textbf{The dynamic-namespace trap.} If the compiler cannot statically know a function's
locals---because it contains \texttt{exec()} against an implicit namespace---it abandons fast
locals and falls back to dictionary-based \texttt{LOAD\_NAME}, searching local, then global,
then builtin namespaces at runtime. (\texttt{from module import *} is itself a
\texttt{SyntaxError} at function scope in Python 3.) Keep function scopes statically
analyzable.

\textbf{\texttt{SyntaxError} vs \texttt{NameError}.} Parsing and symbol analysis happen
before execution, so a malformed expression is rejected at compile time and never reaches the
VM; an unbound name is a runtime \texttt{NameError}. Knowing which stage owns an error tells
you where to look.

\textbf{Don't fight immutability of code objects.} ``Patching bytecode at runtime'' is a
stunt, not an optimization. The supported lever is the \emph{function} wrapper around the
code object---decorators, \texttt{functools.partial}, closures.

\section{Summary and cross-references}

\begin{itemize}
  \item A Python program is a \textbf{graph of heap objects}; names are \textbf{bindings},
    not typed storage. Assignment rebinds and adjusts refcounts; it never copies.
    \texttt{is} is identity, \texttt{==} is value.
  \item Source is \textbf{compiled ahead of time} (tokenizer $\rightarrow$ PEG parser,
    PEP 617 $\rightarrow$ symbol table $\rightarrow$ compiler) into an immutable
    \textbf{\texttt{PyCodeObject}}. Only the \textbf{evaluation loop} is interpretation.
  \item The \textbf{symbol table} statically classifies every name, enabling array-indexed
    \texttt{LOAD\_FAST} locals and \textbf{cell}-based closures.
  \item Modern bytecode (3.11+) includes \texttt{RESUME}, the unified \texttt{BINARY\_OP},
    and fused \textbf{superinstructions}, and it \textbf{specializes itself} as it
    runs---read \texttt{dis} output, not folklore.
  \item Performance is governed by the pipeline: \textbf{locals beat globals}, dynamic
    namespaces defeat optimization, and errors belong to the stage that produces them.
\end{itemize}

\textbf{Cross-references.} The \texttt{PyObject} header and reference counting $\rightarrow$
Chapter 2. Cyclic garbage collection and comprehension/closure scope $\rightarrow$
Chapter 3. The specializing adaptive interpreter in depth $\rightarrow$ Vol VI, Ch 17. The
copy-and-patch JIT (PEP 744) $\rightarrow$ Vol VII. The formal grammar, lexical rules, and
the data model $\rightarrow$ Vol X. The \texttt{ceval.c} evaluation loop in full $\rightarrow$
Vol VIII.
